% Book 4, Year 2, chapter 10 -- Assembly language and the calling convention.
% Provenance (never printed): sources/book4/BRIEFS.md, chapter 10.
\chapter{Assembly Language and the Calling Convention}\label{ch:b2:assembly-and-calling-convention}

What does the C statement \code{x = a[i] + 1;} become inside the processor? Four instructions, each doing one small
thing, traced here one line per instruction as the series' RISC-V machine runs them.

\omcode{code/bachelor-2/10-assembly-and-calling-convention/listings/hook.S}{6}{11}{\code{x = a[i] + 1} in
RISC-V}\label{lst:b2:assembly-and-calling-convention:hook}

\omoutput[1-2,6-9]{code/bachelor-2/10-assembly-and-calling-convention/outputs/hook.txt}{Shift, add, load, add: the
address of \code{a[i]} is computed, then the value fetched.}\label{out:b2:assembly-and-calling-convention:hook}

The C line hid an address computation, a memory access, and the fact that arithmetic happens only between
registers. This chapter reads and writes such code: the instructions of a real instruction set, how they are
encoded in 32-bit words, how C's \code{if}, loops, arrays and structures become them, and the convention that lets
functions compiled separately call one another.

\begin{recall}
Book 2: the toy machine has registers, flags, jumps, and \code{CALL}/\code{RET} through a stack. Book 3: a program
becomes object files, then a process whose memory holds text, data, heap and stack; each call pushes a stack frame;
values are aligned to their size; Book 3 also covers byte order, sign extension and the return address. Two's
complement represents negative integers.
\end{recall}

\section{The machine and its instructions}\label{sec:b2:assembly-and-calling-convention:isa}

\begin{definition}[Instruction set architecture]\label{def:b2:assembly-and-calling-convention:isa}
An \emph{instruction set architecture}\index{instruction set architecture} (ISA) is the contract between a processor
and its programs: the registers, the instructions with their encodings and effects, the memory model, and what
happens on an error. Every processor that implements the ISA runs the same programs, however it is built
(\cref{ch:b2:processor-and-pipelining}).
\end{definition}

This book uses RISC-V, an open ISA whose specification is public, small and regular. The series' machine
\code{rv} implements the base integer set RV32I with the multiply-and-divide extension M, on 32-bit words.

% ledger: b2:rv-origin
\begin{history}[An open instruction set]
RISC-V began in May 2010 as a research project on instruction sets in a university laboratory, meant at first for
teaching and experiments. Its authors published the specification openly; a foundation took it over in 2015, and an
international association has maintained it since 2020. Anyone may build a RISC-V processor without a licence, which
is why it is found today in microcontrollers, storage controllers and research chips.
\end{history}

\begin{definition}[General-purpose register, register file]\label{def:b2:assembly-and-calling-convention:register}
A \emph{general-purpose register}\index{general-purpose register} is a storage cell inside the processor, holding one
word, that instructions name directly as operands. The \emph{register file}\index{register file} is the set of them:
in RV32, 32 registers \code{x0} to \code{x31} of 32 bits, where \code{x0} always reads 0 and ignores writes.
\end{definition}

\begin{definition}[Immediate]\label{def:b2:assembly-and-calling-convention:immediate}
An \emph{immediate}\index{immediate} is a constant operand stored inside the instruction word itself, such as the
\code{1} of \code{addi a2, t1, 1}. In RV32I most immediates have 12 bits, read signed: $-2048$ to $2047$.
\end{definition}

% ledger: b2:rv-abi-names
% ledger: b2:rv-cc-saved
\begin{omfigure}
\centering\footnotesize
\begin{tabular}{@{}llll@{}}
\hline
register & ABI name & role & preserved across calls \\
\hline
\code{x0} & \code{zero} & always 0 & --- \\
\code{x1} & \code{ra} & return address & no \\
\code{x2} & \code{sp} & stack pointer & yes \\
\code{x3}, \code{x4} & \code{gp}, \code{tp} & global and thread pointers & not allocated \\
\code{x5}--\code{x7}, \code{x28}--\code{x31} & \code{t0}--\code{t6} & temporaries & no (caller-saved) \\
\code{x8}, \code{x9}, \code{x18}--\code{x27} & \code{s0}--\code{s11} & saved registers & yes (callee-saved) \\
\code{x10}--\code{x17} & \code{a0}--\code{a7} & arguments; \code{a0}, \code{a1} results & no (caller-saved) \\
\hline
\end{tabular}
\par\smallskip{\small The \omterm{def:b2:assembly-and-calling-convention:register}{register file} under its ABI names, the names programs use. The last column is the
\omterm{def:b2:assembly-and-calling-convention:convention}{calling convention} of \cref{sec:b2:assembly-and-calling-convention:convention}.}
\end{omfigure}

The arithmetic and logic instructions take two registers, or a register and an \omterm{def:b2:assembly-and-calling-convention:immediate}{immediate}, and write a third:
\code{add}, \code{sub}, \code{and}, \code{or}, \code{xor}, the shifts \code{sll}, \code{srl} (logical, zeros in) and
\code{sra} (arithmetic, copies of the sign bit in), the comparisons \code{slt} and \code{sltu} (set to 1 if less,
signed or unsigned), and \code{mul}, \code{div}, \code{rem} from the M extension; each has an \omterm{def:b2:assembly-and-calling-convention:immediate}{immediate} form with an
\code{i} suffix. A first program sums $1$ to $100$:

\omcode{code/bachelor-2/10-assembly-and-calling-convention/listings/first.S}{3}{15}{The sum of the first hundred
integers}\label{lst:b2:assembly-and-calling-convention:first}

\omoutput{code/bachelor-2/10-assembly-and-calling-convention/outputs/first.txt}{$100 \cdot 101 / 2$.}

\begin{definition}[Assembler directive, label, pseudo-instruction]\label{def:b2:assembly-and-calling-convention:directive}
An assembly source line holds at most one statement. An \emph{assembler directive}\index{assembler directive}, a
word starting with a dot (\code{.text}, \code{.data}, \code{.word}, \code{.asciz}), tells the assembler where to put
what follows or which bytes to emit. A \emph{label}\index{label (assembly)}, \code{loop:}, names the address of the
next statement. A \emph{pseudo-instruction}\index{pseudo-instruction} (\code{li}, \code{la}, \code{mv}, \code{j},
\code{ret}, \code{call}) is a convenient name that the assembler expands into one or two real instructions.
\end{definition}

\code{ecall} asks the environment for a service: the number in \code{a7} selects it (1 prints an integer, 11 a
character, 93 exits), \code{a0} carries the argument. \code{li} loads any 32-bit constant; how is the subject of
the encoding section.

\section{Memory}\label{sec:b2:assembly-and-calling-convention:memory}

\begin{definition}[Load--store architecture, base-plus-offset addressing]\label{def:b2:assembly-and-calling-convention:load-store}
In a \emph{load--store architecture}\index{load--store architecture} only loads and stores access memory; every other
instruction works on registers. RISC-V addresses memory by \emph{base plus offset}\index{base-plus-offset
addressing}: \code{lw t1, 8(t0)} loads the word at the address $\code{t0} + 8$, the offset a 12-bit signed \omterm{def:b2:assembly-and-calling-convention:immediate}{immediate}.
\end{definition}

Loads come in three sizes: \code{lw} a word, \code{lh} a half (16 bits) and \code{lb} a byte, the last two
sign-extended to 32 bits, with \code{lhu} and \code{lbu} zero-extending; stores are \code{sw}, \code{sh},
\code{sb}. RISC-V is little-endian, and \code{rv} requires a word access to be at a multiple of 4 and a half at a
multiple of 2.

\omcode{code/bachelor-2/10-assembly-and-calling-convention/listings/memory.S}{2}{20}{Bytes of a word, signed and
unsigned}\label{lst:b2:assembly-and-calling-convention:memory}

\omoutput{code/bachelor-2/10-assembly-and-calling-convention/outputs/memory.txt}{The lowest byte comes first; the
same byte \code{0xff} is $-1$ or 255.}\label{out:b2:assembly-and-calling-convention:memory}

\omoutput{code/bachelor-2/10-assembly-and-calling-convention/outputs/misaligned.txt}{\code{misaligned.S} loads a word
at an address $\equiv 2 \pmod 4$.}\label{out:b2:assembly-and-calling-convention:misaligned}

\begin{omfigure}
\begin{tikzpicture}[font=\scriptsize, r/.style={draw, minimum width=50mm, align=center}]
\node[r, minimum height=11mm, fill=omThm!10] (s) at (0,0) {stack, 8 MiB below \texttt{0x80000000}\\ grows down from \texttt{sp = 0x7ffffff0}};
\node[r, minimum height=8mm] (g) at (0,-0.95) {(not mapped)};
\node[r, minimum height=8mm, fill=omDef!10] (h) at (0,-1.75) {heap, grows up (\texttt{sbrk})};
\node[r, minimum height=6mm, fill=omDef!20] (d) at (0,-2.45) {\texttt{.data}};
\node[r, minimum height=8mm] (g2) at (0,-3.15) {(not mapped)};
\node[r, minimum height=7mm, fill=omProp!15] (t) at (0,-3.9) {\texttt{.text}: instructions, read-only};
\node[r, minimum height=6mm, fill=black!8] (z) at (0,-4.55) {not mapped (a null pointer faults)};
\foreach \n/\a in {s.north west/0x80000000, d.south west/0x10000000, t.south west/0x00010000, z.south west/0x00000000}
  \node[anchor=east, font=\scriptsize\ttfamily] at ([xshift=-2mm]\n) {\a};
\draw[ptr] ([xshift=-12mm]g.north) -- ++(0,-0.3);
\draw[ptr] ([xshift=12mm]g.south) -- ++(0,0.3);
\end{tikzpicture}
\par\smallskip{\small The memory of an \code{rv} program, addresses growing upwards. Every address outside the shaded
regions faults.}
\end{omfigure}

Arrays and strings are walked with a pointer that a loop advances by the element size; a structure is a base
address and fixed offsets, padding included.

\omcode{code/bachelor-2/10-assembly-and-calling-convention/listings/array_sum.S}{6}{17}{A pointer walk: the loop
ends when the pointer reaches the end}\label{lst:b2:assembly-and-calling-convention:array}

\omoutput{code/bachelor-2/10-assembly-and-calling-convention/outputs/array_sum.txt}{The sum of the six page
counts.}

A C string has no length stored anywhere: the code walks bytes until it meets the NUL, and the length is the
difference of two pointers.

\omcode{code/bachelor-2/10-assembly-and-calling-convention/listings/strlen.S}{5}{11}{\code{strlen}: one byte per
turn, \code{lbu} because a byte above 127 is not negative text}\label{lst:b2:assembly-and-calling-convention:strlen}

\omoutput{code/bachelor-2/10-assembly-and-calling-convention/outputs/strlen.txt}{The length of the street
name.}

\omcode{code/bachelor-2/10-assembly-and-calling-convention/listings/struct_access.S}{1}{24}{Fields at offsets 0, 4
and 6 of an 8-byte struct}\label{lst:b2:assembly-and-calling-convention:struct}

\omoutput{code/bachelor-2/10-assembly-and-calling-convention/outputs/struct_access.txt}{The pages of the two books on
shelf B, $65 + 412$.}

\section{Control flow}\label{sec:b2:assembly-and-calling-convention:control}

There are no structured statements in assembly: only conditional branches (\code{beq}, \code{bne}, \code{blt},
\code{bge}, and \code{bltu}, \code{bgeu} for unsigned comparisons), which jump to a \omterm{def:b2:assembly-and-calling-convention:directive}{label} if their two registers
compare as asked, and unconditional jumps, \code{jal} and \code{jalr}. Compiling a statement is a method.

\begin{method}[Compiling an if]\label{met:b2:assembly-and-calling-convention:compile-if}
For \code{if (c) A else B}: evaluate the operands of \code{c}; branch on the \emph{opposite} condition to a \omterm{def:b2:assembly-and-calling-convention:directive}{label}
\code{else}; write \code{A}, then \code{j end}; write \code{else:} and \code{B}; write \code{end:}. Without an
\code{else}, the branch goes directly to \code{end}.
\end{method}

\begin{method}[Compiling a loop]\label{met:b2:assembly-and-calling-convention:compile-loop}
For \code{while (c) B}: write a \omterm{def:b2:assembly-and-calling-convention:directive}{label} \code{test:}, branch on the opposite of \code{c} to \code{out}, write \code{B},
\code{j test}, and \code{out:}. When the body is known to run at least once (\code{do B while (c)}), put the test at
the bottom: \code{loop:}, \code{B}, branch on \code{c} to \code{loop}, one branch per turn instead of two. A
\code{for} loop is its initialisation followed by a \code{while}.
\end{method}

\omcode{code/bachelor-2/10-assembly-and-calling-convention/listings/if_while.S}{3}{27}{An absolute value, and the
$3n + 1$ walk from 27}\label{lst:b2:assembly-and-calling-convention:if-while}

\omoutput{code/bachelor-2/10-assembly-and-calling-convention/outputs/if_while.txt}{$|-42|$, and the 111 steps of the walk from 27 down to 1.}

\begin{omfigure}
\begin{tikzpicture}[font=\scriptsize, b/.style={draw, rounded corners=2pt, fill=omDef!10, align=left, minimum width=30mm}]
\node[b] (t) at (0,0) {\texttt{test:} \texttt{beq t0, t2, out}};
\node[b] (p) at (0,-1.2) {\texttt{andi t1, t0, 1}\\ \texttt{bnez t1, odd}};
\node[b] (e) at (-2.6,-2.6) {\texttt{srli t0, t0, 1}\\ \texttt{j step}};
\node[b] (o) at (2.6,-2.6) {\texttt{odd:} \texttt{slli}, \texttt{add},\\ \texttt{addi}: $3n + 1$};
\node[b] (s) at (0,-3.9) {\texttt{step:} \texttt{addi a0, a0, 1}\\ \texttt{j test}};
\node[b, fill=omThm!10] (x) at (4.6,0) {\texttt{out:} print \texttt{a0}};
\draw[ptr] (t) -- node[left] {$n \ne 1$} (p);
\draw[ptr] (t) -- node[above] {$n = 1$} (x);
\draw[ptr] (p) -- node[left, xshift=-1mm] {even} (e);
\draw[ptr] (p) -- node[right, xshift=1mm] {odd} (o);
\draw[ptr] (e) -- (s); \draw[ptr] (o) -- (s);
\draw[ptr] (s.west) -- ++(-3.6,0) |- (t.west);
\end{tikzpicture}
\par\smallskip{\small The control-flow graph of the \code{while} loop of \cref{lst:b2:assembly-and-calling-convention:if-while}:
each box is a straight run of instructions entered only at its top, each arrow a branch or a fall-through.}
\end{omfigure}

A \code{switch} on a dense range of values compiles to a table of addresses.

\begin{definition}[Jump table]\label{def:b2:assembly-and-calling-convention:jump-table}
A \emph{jump table}\index{jump table} is an array of code addresses, one per case value; a \code{switch} on $v$
checks that $v$ is in range, loads the entry at index $v$, and jumps to it with \code{jr}: one indirect jump,
whatever the number of cases.
\end{definition}

\omcode{code/bachelor-2/10-assembly-and-calling-convention/listings/switch_table.S}{2}{19}{A table of four case
addresses, and the bounds check}\label{lst:b2:assembly-and-calling-convention:switch}

\omoutput{code/bachelor-2/10-assembly-and-calling-convention/outputs/switch_table.txt}{The kinds 1, 0, 3, 7 and 2;
kind 7 fails the bounds check.}\label{out:b2:assembly-and-calling-convention:switch}

\section{Encoding}\label{sec:b2:assembly-and-calling-convention:encoding}

\begin{definition}[Instruction format, opcode]\label{def:b2:assembly-and-calling-convention:format}
Every RV32I instruction is one 32-bit word. Its \emph{opcode}\index{opcode}, the low 7 bits, with the fields
\code{funct3} and \code{funct7}, selects the operation. An \emph{instruction format}\index{instruction format} is
the layout of the other fields; RV32I has six, named by letter.
\end{definition}

% ledger: b2:rv-formats
% ledger: b2:rv-encodings
\begin{omfigure}
\begin{tikzpicture}[font=\tiny, x=3.5mm, y=6mm]
\def\fld#1#2#3#4{\draw[fill=#4] (#1,0) rectangle (#2,1); \node at ({(#1+#2)/2},0.5) {#3};}
\foreach \y/\name in {0/R, -1/I, -2/S, -3/B, -4/U, -5/J} \node[anchor=east, font=\scriptsize\bfseries] at (-0.4,\y+0.5) {\name};
\begin{scope}[yshift=0cm]
\fld{0}{7}{funct7}{black!8}\fld{7}{12}{rs2}{omDef!10}\fld{12}{17}{rs1}{omDef!10}\fld{17}{20}{f3}{black!8}\fld{20}{25}{rd}{omThm!10}\fld{25}{32}{opcode}{black!8}
\end{scope}
\begin{scope}[yshift=-6mm]
\fld{0}{12}{imm[11:0]}{omProp!15}\fld{12}{17}{rs1}{omDef!10}\fld{17}{20}{f3}{black!8}\fld{20}{25}{rd}{omThm!10}\fld{25}{32}{opcode}{black!8}
\end{scope}
\begin{scope}[yshift=-12mm]
\fld{0}{7}{imm[11:5]}{omProp!15}\fld{7}{12}{rs2}{omDef!10}\fld{12}{17}{rs1}{omDef!10}\fld{17}{20}{f3}{black!8}\fld{20}{25}{imm[4:0]}{omProp!15}\fld{25}{32}{opcode}{black!8}
\end{scope}
\begin{scope}[yshift=-18mm]
\fld{0}{7}{imm[12|10:5]}{omProp!15}\fld{7}{12}{rs2}{omDef!10}\fld{12}{17}{rs1}{omDef!10}\fld{17}{20}{f3}{black!8}\fld{20}{25}{imm[4:1|11]}{omProp!15}\fld{25}{32}{opcode}{black!8}
\end{scope}
\begin{scope}[yshift=-24mm]
\fld{0}{20}{imm[31:12]}{omProp!15}\fld{20}{25}{rd}{omThm!10}\fld{25}{32}{opcode}{black!8}
\end{scope}
\begin{scope}[yshift=-30mm]
\fld{0}{20}{imm[20|10:1|11|19:12]}{omProp!15}\fld{20}{25}{rd}{omThm!10}\fld{25}{32}{opcode}{black!8}
\end{scope}
\foreach \b/\x in {31/0.5, 25/6.5, 20/11.5, 15/16.5, 12/19.5, 7/24.5, 0/31.5}
  \node[font=\tiny, text=black!60] at (\x, 1.3) {\b};
\end{tikzpicture}
\par\smallskip{\small The six formats, bit 31 on the left. The register fields stay in the same place in every format,
so the processor can read the registers before it knows the instruction; the \omterm{def:b2:assembly-and-calling-convention:immediate}{immediates} are scattered to make that
possible. B and J store an even offset without its bit 0.}
\end{omfigure}

The assembler's listing shows each word next to its source. Here \code{li} at work:

\omoutput{code/bachelor-2/10-assembly-and-calling-convention/outputs/constants.txt}{\code{constants.S} assembled:
address, word, disassembly, source.}\label{out:b2:assembly-and-calling-convention:constants}

% ledger: b2:rv-lui
\begin{proposition}[Building a 32-bit constant]\label{prop:b2:assembly-and-calling-convention:li}
Let $c$ be a 32-bit word. Let $\mathit{hi} = \lfloor (c + 2^{11}) / 2^{12} \rfloor \bmod 2^{20}$ and let
$\mathit{lo}$ be the low 12 bits of $c$ read as a signed number, in $[-2048, 2047]$. Then \code{lui rd, hi}
followed by \code{addi rd, rd, lo} leaves $c$ in \code{rd} (modulo $2^{32}$).
\end{proposition}

\begin{proof}
\code{lui} places $\mathit{hi} \cdot 2^{12}$ in \code{rd}; \code{addi} adds $\mathit{lo}$, sign-extended. Write $c = 2^{12} q
+ r$ with $0 \le r < 2^{12}$. If $r < 2^{11}$, then $c + 2^{11} = 2^{12} q + (r + 2^{11})$ with $r + 2^{11} < 2^{12}$,
so $\mathit{hi} = q$; bit 11 of $c$ is 0, so $\mathit{lo} = r$, and $2^{12} q + r = c$. If $r \ge 2^{11}$, then $c +
2^{11} = 2^{12} (q + 1) + (r - 2^{11})$, so $\mathit{hi} = q + 1$; bit 11 is 1, so $\mathit{lo} = r - 2^{12}$, negative,
and $2^{12}(q + 1) + r - 2^{12} = c$. All sums are taken modulo $2^{32}$, which \code{lui} and \code{addi} also do.
\end{proof}

The term $2^{11}$ is the rounding: when the low part will be read as negative, the high part must be one larger.
\code{0x12345fff} becomes \code{lui 0x12346} and \code{addi -1} in \cref{out:b2:assembly-and-calling-convention:constants}.
The scattered B and J \omterm{def:b2:assembly-and-calling-convention:immediate}{immediates} also fix how far control can go.

\begin{proposition}[Reach of branches and jumps]\label{prop:b2:assembly-and-calling-convention:ranges}
A conditional branch reaches the addresses $\mathit{pc} - 4096$ to $\mathit{pc} + 4094$, and \code{jal} the
addresses $\mathit{pc} - 2^{20}$ to $\mathit{pc} + 2^{20} - 2$: about $\pm 4$\,KiB and $\pm 1$\,MiB.
\end{proposition}

\begin{proof}
A B \omterm{def:b2:assembly-and-calling-convention:immediate}{immediate} has 12 stored bits, bits 12 to 1 of an even offset, whose bit 0 is 0: a 13-bit two's-complement
number, from $-2^{12}$ to $2^{12} - 2$. A J \omterm{def:b2:assembly-and-calling-convention:immediate}{immediate} stores bits 20 to 1: a 21-bit even number, from $-2^{20}$ to
$2^{20} - 2$. The target is $\mathit{pc}$ plus the offset.
\end{proof}

Decoding reverses the process: read the \omterm{def:b2:assembly-and-calling-convention:format}{opcode}, pick the format, extract the fields. \code{rv dis} does it for a file of
words, as a debugger does for a memory dump:

\omoutput{code/bachelor-2/10-assembly-and-calling-convention/outputs/decode.txt}{Five words of
\code{words.hex}, disassembled.}\label{out:b2:assembly-and-calling-convention:decode}

\section{The calling convention}\label{sec:b2:assembly-and-calling-convention:convention}

A function compiled today must be callable from code compiled years ago, by another compiler. They agree on a
contract.

% ledger: b2:rv-cc-args
\begin{definition}[Calling convention, argument register, return-value register]\label{def:b2:assembly-and-calling-convention:convention}
A \emph{calling convention}\index{calling convention} fixes how functions pass arguments and results, which
registers a call may change, and how the stack is used. In RISC-V the integer arguments go in the \emph{argument
registers}\index{argument register} \code{a0} to \code{a7}, in order (more on the stack), and the result comes back in
the \emph{return-value registers}\index{return-value register} \code{a0} and \code{a1}. \code{jal ra, f}
(\code{call f}) jumps and leaves the return address in \code{ra}; \code{jalr zero, 0(ra)} (\code{ret}) goes back.
\end{definition}

\begin{definition}[Caller-saved register, callee-saved register]\label{def:b2:assembly-and-calling-convention:saved}
A \emph{caller-saved register}\index{caller-saved register} (\code{t0}--\code{t6}, \code{a0}--\code{a7},
\code{ra}) may be changed by any call: a caller that needs its value after the call must save it first. A
\emph{callee-saved register}\index{callee-saved register} (\code{s0}--\code{s11}, \code{sp}) must have the same
value after the call as before: a function that uses one saves it on entry and restores it before returning.
\end{definition}

\begin{definition}[Leaf function, prologue, epilogue]\label{def:b2:assembly-and-calling-convention:prologue}
A \emph{leaf function}\index{leaf function} calls no other function. A non-leaf function, or one that uses
\omterm{def:b2:assembly-and-calling-convention:saved}{callee-saved registers}, opens with a \emph{prologue}\index{prologue}, which allocates its stack frame and saves
\code{ra} and the \omterm{def:b2:assembly-and-calling-convention:saved}{callee-saved registers} it will use, and ends with an \emph{epilogue}\index{epilogue}, which
restores them, frees the frame and returns.
\end{definition}

% ledger: b2:rv-sp-align
\begin{method}[Writing a prologue and an epilogue]\label{met:b2:assembly-and-calling-convention:frame}
\begin{enumerate}
\item List what must survive a call: \code{ra} if the function calls; each value used after a call, to be kept in
an \code{s} register.
\item Frame size: 4 bytes per saved register plus locals, rounded up to a multiple of 16 (the stack pointer stays
16-byte aligned at every call).
\item \omterm{def:b2:assembly-and-calling-convention:prologue}{Prologue}: \code{addi sp, sp, -size}; \code{sw ra, size-4(sp)}; \code{sw s0, size-8(sp)}, and so on.
\item \omterm{def:b2:assembly-and-calling-convention:prologue}{Epilogue}, in reverse: the loads, then \code{addi sp, sp, size}, then \code{ret}. Every path out of the
function goes through it.
\end{enumerate}
\end{method}

\omcode{code/bachelor-2/10-assembly-and-calling-convention/listings/factorial.S}{14}{28}{A recursive function under the
convention}\label{lst:b2:assembly-and-calling-convention:fact}

\omoutput[1-2,5-13,26-30,33-39]{code/bachelor-2/10-assembly-and-calling-convention/outputs/factorial.txt}{The trace
of \code{fact(3)}: the first call, the deepest one, and the first return.}\label{out:b2:assembly-and-calling-convention:fact}

\begin{omfigure}
\begin{tikzpicture}[font=\scriptsize, w/.style={draw, minimum width=34mm, anchor=north west, inner sep=1pt}]
\foreach \i/\who/\ra/\so/\a in {0/main/{ra = 0xffff0000}/{(unused)}/0x7fffffe0,
    1/{fact(3)}/{ra = 0x00010010}/{s0 = 0 (main's)}/0x7fffffd0, 2/{fact(2)}/{ra = 0x0001004c}/{s0 = 3}/0x7fffffc0,
    3/{fact(1)}/{ra = 0x0001004c}/{s0 = 2}/0x7fffffb0, 4/{fact(0)}/{ra = 0x0001004c}/{s0 = 1}/0x7fffffa0} {
  \node[w, minimum height=4.5mm, fill=omDef!10] at (0,-1.3*\i) {\texttt{\ra}};
  \node[w, minimum height=4.5mm, fill=omThm!10] at (0,-1.3*\i-0.45) {\texttt{\so}};
  \node[w, minimum height=4mm, text=black!50] at (0,-1.3*\i-0.9) {8 bytes unused};
  \node[anchor=west] at (4.3,-1.3*\i-0.65) {\who};
  \node[anchor=east, font=\scriptsize\ttfamily] at (-0.15,-1.3*\i-1.3) {\a};
  \node[anchor=west, font=\tiny, text=black!60] at (3.5,-1.3*\i-0.22) {+12};
  \node[anchor=west, font=\tiny, text=black!60] at (3.5,-1.3*\i-0.67) {+8};
}
\draw[ptr] (6.0,-6.5) node[anchor=west, text=black!70] {\texttt{sp} at step 28} -- (3.5,-6.5);
\end{tikzpicture}
\par\smallskip{\small The stack when \code{fact(0)} has saved its registers (step 28), drawn from the trace: five
16-byte frames, each with the saved \code{ra} at offset 12 and the saved \code{s0} at offset 8 (the other 8 bytes
unused). Each \code{fact(n)} saved its caller's \code{n}, which it restores on the way back.}
\end{omfigure}

The convention is a promise each function makes to its callers; breaking it breaks code far away. Here a helper
uses \code{s0} without saving it:

\omcode{code/bachelor-2/10-assembly-and-calling-convention/listings/clobber_bug.S}{5}{30}{A callee that clobbers
\code{s0}}\label{lst:b2:assembly-and-calling-convention:clobber}

\omoutput{code/bachelor-2/10-assembly-and-calling-convention/outputs/clobber_bug.txt}{61 instead of
130.}\label{out:b2:assembly-and-calling-convention:clobber}

\code{fine} sets \code{s0} to 21 at each call, so the caller's running total restarts from 21: $21 + 0$, then $21 +
40 = 61$. Nothing in either function alone is wrong for its own purpose; the bug is the broken contract, which is why
it is so hard to find in a large program. \code{fine} should use \code{t0}, which it may change freely.

\section{Reading a compiler's output}\label{sec:b2:assembly-and-calling-convention:compiler}

A compiler does the same work. \code{gcc -O1 -S} writes the assembly of a C file for the machine at hand, here
x86-64, an older and larger ISA with 16 registers, two-operand instructions, and memory operands inside arithmetic.

\omoutput[3-21]{code/bachelor-2/10-assembly-and-calling-convention/outputs/compiled.txt}{\code{sum(a, n)} compiled
by gcc for x86-64, directives removed (\code{compiled.sh}).}\label{out:b2:assembly-and-calling-convention:gcc-sum}

% ledger: b2:x86-64-args
The ideas are the same. The x86-64 convention passes the first integer arguments in \code{rdi}, \code{rsi},
\code{rdx}, \code{rcx}, \code{r8}, \code{r9} and returns in \code{rax} (\code{eax} is its low half): here \code{a} in
\code{rdi}, \code{n} in \code{esi}. The loop is a pointer walk like \cref{lst:b2:assembly-and-calling-convention:array},
with the end pointer computed once (\code{leaq}) and the test at the bottom; the case $n \le 0$ is handled by a
branch around the loop. The first instruction, \code{endbr64}, only marks a legal target for indirect jumps, a
protection of recent x86-64 processors. And the \code{switch} of five cases became a \omterm{def:b2:assembly-and-calling-convention:jump-table}{jump table}:

\omoutput[22-36]{code/bachelor-2/10-assembly-and-calling-convention/outputs/compiled.txt}{\code{fine_per_day}: one
unsigned comparison (\code{ja}), then an indirect jump through \code{.L9}.}\label{out:b2:assembly-and-calling-convention:gcc-switch}

\begin{inpractice}[Reading assembly]
Read compiler output to check what the optimiser did, not to write it: compile with \code{-O1} or \code{-O2} and
\code{-S}, or use \code{objdump -d} on an object file. Find the arguments in their registers first, then the loop
(a backward branch), then the stores. When a program fails under a debugger, the disassembly and the registers are
the ground truth; the convention tells you which values the caller can still trust.
\end{inpractice}

\section{Exercises}\label{sec:b2:assembly-and-calling-convention:exercises}

\begin{exercise}[$\star$]\label{exo:b2:assembly-and-calling-convention:1}
Predict, in decimal, the values of \code{a1} to \code{a5} after these instructions, then check with the register
dump in the solution.
\omcode{code/bachelor-2/10-assembly-and-calling-convention/listings/exo_regs.S}{3}{8}{Shifts, comparisons and
\code{xori}}
\end{exercise}

\begin{exercise}[$\star$]\label{exo:b2:assembly-and-calling-convention:2}
This program prints 1. Explain why with \cref{met:b2:assembly-and-calling-convention:compile-loop}, and fix the
branch.
\omcode{code/bachelor-2/10-assembly-and-calling-convention/listings/exo_bge.S}{2}{10}{A loop that runs once}
\end{exercise}

\begin{exercise}[$\star$]\label{exo:b2:assembly-and-calling-convention:3}
Write \code{max(a0 = address, a1 = n)}, a leaf that returns the largest of $n \ge 1$ signed words, and test it on
the temperatures $-12, 4, -3, 19, 7, -40$. Which branch would be wrong for unsigned words?
\end{exercise}

\begin{exercise}[$\star$]\label{exo:b2:assembly-and-calling-convention:4}
Encode by hand, in hexadecimal, using the formats of this chapter: \code{addi a0, a0, -1}; \code{sw ra, 12(sp)};
and \code{beq a0, zero, main} placed 8 bytes after \code{main}. (\code{a0} is \code{x10}, \code{ra} \code{x1},
\code{sp} \code{x2}; the \omterm{def:b2:assembly-and-calling-convention:format}{opcodes} are \code{0010011}, \code{0100011} and \code{1100011}; \code{funct3} is 0, 2 and 0.)
\end{exercise}

\begin{exercise}[$\star\star$]\label{exo:b2:assembly-and-calling-convention:5}
Decode by hand the five words \code{00b50633}, \code{00452283}, \code{fe0518e3}, \code{00c000ef} and
\code{0005a503}, as if they started at address \code{0x00010000}. For the branch and the jump, give the target.
\end{exercise}

\begin{exercise}[$\star\star$]\label{exo:b2:assembly-and-calling-convention:6}
This program never ends; \code{rv} stops it at the step limit. Explain, following \code{ra}, where each \code{ret}
goes, and fix \code{double_plus_two}.
\omcode{code/bachelor-2/10-assembly-and-calling-convention/listings/exo_ra.S}{12}{18}{A non-leaf function}
\omoutput[2]{code/bachelor-2/10-assembly-and-calling-convention/outputs/exo_ra.txt}{The step limit.}
\end{exercise}

\begin{exercise}[$\star\star$]\label{exo:b2:assembly-and-calling-convention:7}
Write \code{strcmp(a0, a1)} under the convention: negative, zero or positive as the first string is smaller, equal or
larger, comparing bytes as unsigned. Test it on \code{Mill Lane} against \code{Mill Row}, a string against itself,
and \code{Mill} against \code{Mill Lane}.
\end{exercise}

\begin{exercise}[$\star\star$]\label{exo:b2:assembly-and-calling-convention:8}
(a) Show that the pair $(\mathit{hi}, \mathit{lo})$ of \cref{prop:b2:assembly-and-calling-convention:li} is the only
one with $\mathit{lo} \in [-2048, 2047]$ and $\mathit{hi} \in [0, 2^{20})$. (b) For which constants does \code{li}
need only one instruction, \code{addi} or \code{lui}? (c) Compute $\mathit{hi}$ for \code{0x7ffff800} and check,
modulo $2^{32}$, that the pair still works.
\end{exercise}

\begin{exercise}[$\star\star$]\label{exo:b2:assembly-and-calling-convention:9}
Compile by hand \code{int fine(const struct loan *l)}, where \code{struct loan { int days; short rate; char lost; }},
which returns 2000 if \code{lost}, else $(\mathit{days} - 21) \cdot \mathit{rate}$ when positive, else 0. Give the
offsets and the size of the structure first.
\end{exercise}

\begin{exercise}[$\star\star\star$]\label{exo:b2:assembly-and-calling-convention:10}
Write a recursive binary search \code{find(a0 = array, a1 = lo, a2 = hi, a3 = key)} that returns the index of
\code{key} in the sorted range \code{a[lo..hi)}, or $-1$, under the convention, and test it on the seven sorted page
counts for keys present, absent, smaller than all and larger than all. Which registers must it save, and why are its
recursive calls \omterm{def:b2:functional-programming:tail}{tail calls} in disguise?
\end{exercise}

\begin{exercise}[$\star\star\star$]\label{exo:b2:assembly-and-calling-convention:11}
A naive recursive \code{fib(n)} saves \code{ra}, \code{s0} and \code{s1} in a 16-byte frame, and returns at once,
with no frame, for $n < 2$. (a) Prove that it makes $C(n) = 2F_{n+1} - 1$ calls, where $F$ is the Fibonacci sequence.
(b) What is the deepest stack, in frames and in bytes? (c) The test of the solution counts the \code{jal} and
\code{jalr} executed: what number should it expect?
\end{exercise}

\begin{exercise}[$\star\star\star$]\label{exo:b2:assembly-and-calling-convention:12}
Compile with a \omterm{def:b2:assembly-and-calling-convention:jump-table}{jump table} \code{switch (shelf) { case 3: 'A'; case 4: case 5: 'B'; case 7: 'D'; default: '?' }}.
Use one unsigned comparison for the bounds check, and explain why it also rejects values below 3 and negative ones.
\end{exercise}

\section{Problem: Hand-compiling the library's check digit}\label{sec:b2:assembly-and-calling-convention:problem}

\begin{problem}[{Weekend problem --- Hand-compiling the library's check digit}]\label{pb:b2:assembly-and-calling-convention:1}
The library's book codes end with a check digit computed from the other digits, so that a mistyped digit is caught at
the desk. This problem compiles the function by hand, calls a helper under the convention, and proves exactly how many
instructions it executes.

\medskip\noindent\textbf{Part I --- The C function.}\par\nopagebreak
\omcode{code/bachelor-2/10-assembly-and-calling-convention/listings/check_digit.c}{4}{18}{The check digit in C}
\begin{enumerate}[beginpenalty=10000]
\item Compute by hand the check digit of 4006, and compare with the output in the solution.
\item Why a \code{do}--\code{while} rather than a \code{while}? What would \code{check_digit(0)} return?
\item Show that changing one digit of a code always changes the check digit. (Consider the weights 1 and 3
modulo 10.)
\item Swapping two adjacent digits $x$ and $y$ changes the weighted sum by $\pm 2(x - y)$. When is the swap not
caught?
\item Which values live across the call to \code{weigh}? Which may be lost?
\item How large a code fits in a 32-bit unsigned word, in digits?
\end{enumerate}

\medskip\noindent\textbf{Part II --- Compiled by hand.}\par\nopagebreak
\omcode{code/bachelor-2/10-assembly-and-calling-convention/listings/count.S}{5}{22}{The \omterm{def:b2:assembly-and-calling-convention:prologue}{prologue} and the loop}
\begin{enumerate}[resume, beginpenalty=10000]
\item Which registers hold \code{code}, \code{i} and \code{sum}, and why \code{s} registers?
\item The frame is 16 bytes for four saved registers. Write the \omterm{def:b2:assembly-and-calling-convention:prologue}{epilogue}, and check it against \code{count.S}.
\item Which instructions compute \code{code} modulo 10 and \code{code / 10}? Why the unsigned forms?
\item Where is the loop's test, and which method of this chapter was applied?
\item Rewrite the loop to keep \code{sum} in \code{a2} instead of \code{s2}. What breaks?
\item The final expression uses \code{remu} twice. Could the first be avoided?
\end{enumerate}

\medskip\noindent\textbf{Part III --- A helper under the convention.}\par\nopagebreak
\omcode{code/bachelor-2/10-assembly-and-calling-convention/listings/count.S}{40}{46}{\code{weigh}, a leaf}
\begin{enumerate}[resume, beginpenalty=10000]
\item Why does \code{weigh} need no \omterm{def:b2:assembly-and-calling-convention:prologue}{prologue} at all?
\item \code{weigh} computes the weight without a branch. Prove that \code{3 - 2 * (i & 1)} is 3 for even $i$ and 1
for odd $i$.
\item \code{weigh} changes \code{t0} and \code{t1}. Is that allowed? Which caller would it break, and how is that
caller's bug avoided?
\item Inline \code{weigh} into the loop. How many instructions per digit are saved?
\item A real assembler emits \code{call} as two instructions, \code{auipc} and \code{jalr}, for a target that may be far.
How would that change the count?
\item The chapter's test also runs the codes 12, 123456789 and 100000000. Predict their instruction counts.
\end{enumerate}

\medskip\noindent\textbf{Part IV --- Every instruction counted.}\par\nopagebreak
\omoutput{code/bachelor-2/10-assembly-and-calling-convention/outputs/run_count.txt}{\code{rv --stats} on codes of
1, 4 and 9 digits.}
\begin{enumerate}[resume, beginpenalty=10000]
\item Count the instructions executed per digit, from \code{digit:} through \code{bnez}, including \code{weigh}.
\item Count the instructions executed once, before and after the loop.
\item Prove that a code of $n$ digits executes exactly $an + b$ instructions, and give $a$ and $b$.
\item Check the formula against the three counts of the output.
\item Why does the count depend only on $n$, not on the digits themselves?
\item State the result of the problem.
\end{enumerate}
\end{problem}
